\documentclass[a4paper,11pt]{article}

% Encodage et langue
\usepackage[T1]{fontenc}
\usepackage[utf8]{inputenc}
\usepackage[french]{babel}
\usepackage{csquotes}

% Typographie et mise en page
\usepackage{microtype}
\usepackage[a4paper, margin=2.5cm]{geometry}

% Math
\usepackage{amsmath,amssymb,amsthm}

% Graphiques (au cas où)
\usepackage{graphicx}

% Liens et références
\usepackage[colorlinks=true, linkcolor=blue!50!black, citecolor=blue!50!black, urlcolor=blue!50!black]{hyperref}
\usepackage[capitalize,noabbrev]{cleveref}

% Commandes personnalisées (exemples)
\newcommand{\N}{\mathbb{N}}
\newcommand{\Z}{\mathbb{Z}}
% commentaire ponctuel ajouté pour recompiler.

\title{Les nombres premiers de Mersenne}
\author{Paul}
\date{\today}

\begin{document}

\maketitle

\begin{abstract}
Cet article court présente la définition des nombres de Mersenne et donne deux exemples classiques. Nous définissons $M_p$ par une formule simple et calculons $M_3$ et $M_7$.
\end{abstract}

\section{D\'efinition}
Les nombres de Mersenne sont les entiers de la forme
\[
M_p = 2^{p} - 1,
\]
o\`u $p$ est un entier positif. On appelle "nombre de Mersenne premier" un nombre de Mersenne qui est lui-m\^eme premier. Il est clair que si $M_p$ est premier, alors $p$ doit \^etre premier; la r\'eciproque n'est pas toujours vraie.

\section{Exemples: $M_3$ et $M_7$}
Calculons deux exemples simples :
\begin{align*}
M_3 &= 2^{3} - 1 = 8 - 1 = 7, \\%
M_3 &= 7 \quad(\text{qui est premier}),
\\
M_7 &= 2^{7} - 1 = 128 - 1 = 127, \\%
M_7 &= 127 \quad(\text{qui est aussi premier}).
\end{align*}
Ces deux valeurs sont des exemples historiques de nombres de Mersenne premiers.

\end{document}
